\documentclass[11pt,letterpaper]{article}
\usepackage[T1]{fontenc}
\usepackage[utf8]{inputenc}
\usepackage{amsmath,amsthm,mathtools}
\usepackage{newtxtext,newtxmath}
\usepackage[margin=0.88in,headheight=15pt]{geometry}
\usepackage{microtype}
\usepackage{booktabs,longtable,tabularx,array}
\usepackage{xcolor,listings,enumitem,fancyhdr,xurl}
\usepackage[most]{tcolorbox}
\usepackage{titlesec}
\definecolor{ink}{HTML}{3D5044}
\definecolor{soft}{HTML}{F3F4EF}
\definecolor{rulegray}{HTML}{CAD0C4}
\usepackage[colorlinks=true,linkcolor=ink,citecolor=ink,urlcolor=ink]{hyperref}
\hypersetup{pdftitle={AsymptoticAnalysis: a nonduplicative audit of defining-sum integration},
 pdfauthor={Technical review prepared for Vladimir Reshetnikov},
 pdfsubject={Wolfram 15 observations, Mathics source review, Dirichlet dispatch and term-goal fixes}}
\titleformat{\section}{\large\sffamily\bfseries\color{ink}}{\thesection}{0.7em}{}
\titleformat{\subsection}{\normalsize\sffamily\bfseries\color{ink}}{\thesubsection}{0.7em}{}
\titleformat{\subsubsection}{\normalsize\sffamily\bfseries}{\thesubsubsection}{0.7em}{}
\lstset{basicstyle=\ttfamily\footnotesize,breaklines=true,columns=fullflexible,
 keepspaces=true,showstringspaces=false,backgroundcolor=\color{soft},
 frame=single,rulecolor=\color{rulegray},framesep=6pt,
 aboveskip=10pt,belowskip=10pt,tabsize=2}
\setlength{\parindent}{0pt}
\setlength{\parskip}{0.45em}
\setlength{\emergencystretch}{3em}
\renewcommand{\arraystretch}{1.15}
\setlength{\tabcolsep}{4pt}
\setlist{nosep,leftmargin=1.5em}
\newcolumntype{P}[1]{>{\raggedright\arraybackslash}p{#1}}
\newcommand{\code}[1]{\nolinkurl{#1}}
\newcommand{\OO}{\mathcal O}
\newcommand{\rev}{8cee870994f506b501bae3ea6bd4a3a7edb895c1}
\newcommand{\src}[2]{\href{https://github.com/VladimirReshetnikov/Asymptotic/blob/8cee870994f506b501bae3ea6bd4a3a7edb895c1/#1}{#2}}
\newtheorem{proposition}{Proposition}[section]
\newtheorem{lemma}[proposition]{Lemma}
\newtcolorbox{finding}{colback=soft,colframe=rulegray,boxrule=0.5pt,
 arc=0pt,left=9pt,right=9pt,top=8pt,bottom=8pt,breakable}
\pagestyle{fancy}\fancyhf{}
\fancyhead[L]{\small\textsc{AsymptoticAnalysis: incremental technical audit}}
\fancyhead[R]{\small\texttt{8cee870994f5}}
\fancyfoot[L]{\small 10 September 2026}
\fancyfoot[R]{\thepage}
\begin{document}
\begin{titlepage}
{\Large\sffamily\color{ink}\textsc{Source-pinned technical review}\par}
\vspace{0.7in}
{\Huge\bfseries\color{ink}AsymptoticAnalysis\par}
\vspace{0.2in}
{\LARGE Defining-sum integration,\\request semantics, and\\an integer-indexed Dirichlet algebra\par}
\vspace{0.35in}
{\large New findings after the existing review corpus;\\focused repairs, mathematical contracts, and executable artifacts\par}
\vspace{0.45in}
Prepared for Vladimir Reshetnikov\\
10 September 2026
\vfill
\textbf{Repository}\par
\url{https://github.com/VladimirReshetnikov/Asymptotic}\par
\textbf{Reviewed commit}\par
\texttt{\rev}\par
\medskip\hrule\medskip
\textbf{Evidence in this review.} Selected public results and a three-edit candidate
were executed in Wolfram Language 15.0.0. The independent Python artifacts pass
32 test methods. Mathics compatibility was examined in source, but no Mathics
runtime was available. Neither the full upstream suite nor the complete packaged
Wolfram probe file was run. The two new findings are coverage and request-contract
issues; neither is presented as a newly discovered false asymptotic theorem.
\end{titlepage}

\begin{abstract}
This audit examined the canonical forward and inverse machinery, series
operations, exact interval primitives, several special-function modules, and
Mathics compatibility layers, against the repository's extensive existing
review record. Most plausible candidates encountered in those areas were
already documented, already guarded, or not substantiated. They are not
reintroduced as findings.

Two reproducible additions remain in the defining-sum adapter. First, the
package constructor accepts bare large-argument Zeta and Lerch expressions,
but rejects simple affine translations of those expressions even though its
existing series arithmetic can construct the corresponding results. Second,
the same adapter ignores an explicitly supplied \code{SeriesTermGoal} whenever
an explicit cutoff is present. Its result can record a requested goal of one
and a returned count of three. In the Zeta path this also causes a resource
refusal for a request whose term goal makes the expensive work unnecessary.

The article isolates both mechanisms, supplies a narrowly scoped term-goal
patch and a restricted affine-lifting prototype, and develops an additional
integer-indexed Dirichlet-jet prototype. The latter implements exact coefficient
addition, convolution, reciprocal prefixes, and elementary global coefficient
majorants. Its purpose is a concrete development option for this particular
family, rather than another proposal for an unrestricted transseries framework.
Native observations, independent tests, candidate mechanisms, and unexecuted
artifacts are recorded separately throughout.
\end{abstract}
\setcounter{tocdepth}{1}
\hypersetup{bookmarksdepth=2}
\tableofcontents
\clearpage

\section{Assessment, scope, and nonduplication}
\subsection{What should change first}
The smallest useful upstream change is to make the defining-sum adapter obey
both explicit request constraints. It requires three local edits, preserves
the existing coefficient and remainder formulas, and has successful focused
Wolfram controls. The next change should make rational affine wrappers around
an admitted defining-sum atom reach that atom's constructor, rather than falling
through to a native coefficient importer that cannot interpret it.

These are deliberately modest priorities. The retained review corpus already
contains substantial correctness, performance, domain, precision, and
architecture recommendations. Recommending those same projects again would
obscure the incremental contribution requested here.~\cite{reviewindex,register}

\begin{longtable}{@{}P{0.07\textwidth}P{0.15\textwidth}P{0.49\textwidth}P{0.22\textwidth}@{}}
\toprule
ID & Classification & Finding & Evidence\\\midrule\endhead
N01 & Medium; coverage and UX & Bare Zeta/Lerch atoms are admitted, but their simple affine translations are refused by the public Package constructor. Existing arithmetic works around the dispatch gap. & Four native constructor observations; two existing-arithmetic controls.\\
N02 & Medium; request semantics and work selection & Explicit cutoffs disable the term goal in both defining-sum constructors. Zeta can reject a cheap one-term request on the basis of work it need not perform. & Native count mismatches and resource refusal; five candidate controls.\\
D01 & Development option, not a defect & Integer-indexed Dirichlet jets with exact convolution and reciprocal coefficient bounds. & Implemented independent Python prototype; mathematical derivation and tests.\\
\bottomrule
\end{longtable}

The severity labels concern ordinary correctness of the public request contract
and usefulness of the API. N01 returns a failure rather than a wrong expression.
N02 returns a mathematically valid longer expansion, but does not honor the
requested term count. No security severity or comprehensive acceptance claim
is implied.

\subsection{Snapshot and retrieval boundary}
The review is pinned to commit \texttt{\rev}, not to an unspecified moving
\code{main}. The root standalone was imported as text and then loaded from a
local temporary file in the Wolfram evaluator. The import contained 692,380
characters. A smoke call for \code{Sin[x]} at cutoff three returned a
\code{GeneralizedSeries} with normal expression $x$, consistent with the
package's exclusive cutoff convention.~\cite{core}

The substantive source inspection covered 21 canonical modules to different
depths. These include the main kernel, ordinary and composite arithmetic,
series operations, Dirichlet special functions, Gamma/Barnes forward paths,
selected inverse operations, Fourier coefficients, interval certificates, and
six Mathics compatibility modules. The coverage record in
\code{evidence/source-coverage.json} distinguishes full returned modules from
selected or overlapping ranges. This is substantial targeted analysis, not a
claim that every source line, test, script, or historical document received an
equal audit.

\subsection{How earlier findings were excluded}
At the pinned snapshot, the review index contains 35 retained packages across
five waves, plus ten retirement tombstones. The consolidated status register,
all five wave indexes, and selected original discussions were used as the
primary suppression baseline. A further keyword-and-context scan read 62 TeX
files from the retained articles, including split sections and standalone
copies. The scan for waves one through three explicitly reported 47 reads and
no unread files; the remaining fifteen files belong to waves four and five.
This is not a count of 62 independent reviews.~\cite{reviewindex,wave3,wave5,register}

The lexical scan looked for Dirichlet, Zeta, Lerch, divisor, and multiplicative
algebra terminology and examined the returned contexts. It supplements, rather
than replaces, semantic comparison with the maintained finding summaries. It
cannot prove the absence of an equivalent idea expressed with different words.
The new entries below therefore include their nearest related earlier entries
and explain the difference. The machine-readable ledger is
\code{evidence/novelty-ledger.json}.

One public wrong-result candidate encountered during the investigation overlapped
an earlier sided-germ obligation. It was suppressed rather than used to inflate
the new-finding count. The same rule was applied to previously recorded
certificate, truncation, native-boundary, and general budget topics. Their
reproduction details and recommendations are not repeated here.

\subsection{Execution evidence and its limits}
\begin{longtable}{@{}P{0.25\textwidth}P{0.39\textwidth}P{0.29\textwidth}@{}}
\toprule
Evidence population & What was obtained & What it does not establish\\\midrule\endhead
Wolfram baseline & Public constructor, term-count, resource-refusal, and arithmetic observations on version 15.0.0 for Linux x86-64, build dated 6 May 2026. & Full upstream-suite acceptance or equivalent behavior on every supported platform.\\
Term-goal candidate & Three guarded edits applied to the imported standalone; five focused outputs including unchanged no-goal controls. & Full rebuilt modular distribution acceptance, combined changes, or Mathics execution.\\
Independent Python & 28 arithmetic/majorant test methods and four synthetic patch-anchor fixture methods; all 32 passed. & Execution of the Wolfram package or proof of an arbitrary caller-declared infinite coefficient bound.\\
Affine prototype & Packaged standalone helper and a probe driver; the underlying existing arithmetic operations succeeded in separate native calls. & Successful execution of the complete helper file. Its grouped mechanism probe failed at the connector boundary.\\
Mathics & Source-level examination of the compatibility layers and the vocabulary used by the proposed changes. & A Mathics baseline, patched run, benchmark, or compatibility certification.\\
\bottomrule
\end{longtable}

Several requests to the remote evaluator failed with transport errors. A direct
remote \code{Get} also received a truncated source stream; importing the complete
text and loading a local file succeeded. These events are recorded as execution
infrastructure limitations, not repository syntax or algorithm defects.

\section{Where the new issues occur}
\subsection{Two special-function inputs, two natural coordinates}
The defining-sum module supplies constructions unavailable to a straightforward
ordinary Taylor path. For large real $S$, the Zeta adapter uses the convergent
sum
\begin{equation}
 \zeta(S)=\sum_{n=1}^{\infty} n^{-S},\qquad S>1,
\end{equation}
and chooses $w=e^{-S}$. Its exact weights are $\log n$, since
$n^{-S}=w^{\log n}$. An exclusive cutoff $h$ retains precisely the terms with
$\log n<h$. The $n=1$ term is the constant one and is counted as a term by this
adapter.~\cite{dirichlet,dlmf-zeta}

The Lerch adapter instead treats fixed real parameters $z,s$ with $|z|<1$ and
an argument $A$ tending to positive infinity. It uses $w=1/A$, with blocks of
weights $s+k$ and coefficients
\begin{equation}\label{eq:lerch-coeff}
 c_k=\frac{(-1)^k(s)_k}{k!}M_k(z),\qquad
 M_0(z)=\frac1{1-z},\quad M_k(z)=\operatorname{Li}_{-k}(z)\ (k>0).
\end{equation}
The defining series is $\Phi(z,s,A)=\sum_{n\geq0}z^n(A+n)^{-s}$ in the admitted
real domain. For the witness $z=1/2,s=1$, the first coefficients are
$2,-2,6,-26,\ldots$.~\cite{dirichlet,dlmf-lerch}

The original module already attaches explicit remainder metadata to these
constructions. This article does not re-present the module's existing tail
bounds as a new result or a newly discovered need for majorants. The findings
are in admission and request selection around that coefficient machinery.

\subsection{The constructor and the arithmetic layer have different entry points}
The important boundary is not between a ``capable Wolfram'' and an ``incapable
Mathics.'' The source contains a dedicated dispatcher,
\code{dirichletSpecialForwardExpansion}, whose initial guard admits only a
whole expression matching
\begin{lstlisting}
Zeta[_] | LerchPhi[_, _, _]
\end{lstlisting}
An affine wrapper has head \code{Plus} or \code{Times}, so this guard declines it.
Subsequent generic routes do not automatically convert its special-function
leaf into the result that the defining-sum adapter would have returned.
That is the mechanism behind N01.~\cite{dirichlet}

The ordinary arithmetic layer is a separate consumer of already constructed
\code{GeneralizedSeries} objects. It has guarded routes for regular operands,
compatible ordered series, and composite envelopes. Consequently, the fact
that a wrapped expression fails during construction does not imply that its
sum cannot be represented after the leaf is constructed.~\cite{arithmetic,envelope,operations}

This distinction suggests a local integration change: make a small admitted
expression grammar construct its special leaf and then reuse the existing
arithmetic. It does not require changing the mathematical definition of
Zeta, modifying native \code{Series}, or admitting all expressions containing
special functions.

\section{N01: an admitted atom becomes unsupported when translated}
\begin{finding}
\textbf{New coverage/UX finding; successful Wolfram baseline observations.}\par
\textbf{Location:} the whole-expression guard in
\code{dirichletSpecialForwardExpansion}, in
\code{src/Kernel/DirichletSpecialFunctions.wl}.\par
\textbf{Boundary:} explicit \code{"Backend" -> "Package"}. Automatic/native
behavior outside the observed calls is not inferred.
\end{finding}

\subsection{Public witnesses}
With the pinned package loaded, the following call succeeds:
\begin{lstlisting}
AsymptoticAnalysis`AsymptoticExpansion[
  Zeta[x], {x, Infinity, 2}, "Backend" -> "Package"]
\end{lstlisting}
Its normal expression is
\begin{equation}
 1+2^{-x}+3^{-x}+4^{-x}+5^{-x}+6^{-x}+7^{-x},
\end{equation}
a nonzero remainder being recorded as
\code{PowerLogRemainder[E^(-x),Log[8],0]}.
The translated input
\begin{lstlisting}
AsymptoticAnalysis`AsymptoticExpansion[
  Zeta[x] - 1, {x, Infinity, 2}, "Backend" -> "Package"]
\end{lstlisting}
returns \code{Failure["UnsupportedNativeCoefficient",...]} instead. The
reported unsupported coefficient contains a retained
\code{Zeta[u$12^(-1)]}; the generated local symbol is incidental. The message
asks for polynomial logarithmic coefficients and bounded real sine/cosine
modes. It does not tell the caller that the nearby bare expression has a
successful defining-sum constructor.

The Lerch pair is equally simple:
\begin{lstlisting}
AsymptoticAnalysis`AsymptoticExpansion[
  LerchPhi[1/2, 1, x], {x, Infinity, 2},
  "Backend" -> "Package"]
(* Normal: 2/x; remainder: O(x^-2) *)

AsymptoticAnalysis`AsymptoticExpansion[
  1 + LerchPhi[1/2, 1, x], {x, Infinity, 2},
  "Backend" -> "Package"]
(* Failure["UnsupportedNativeCoefficient", ...] *)
\end{lstlisting}
Both successful and failed outputs are transcribed in
\code{evidence/native-observations.json}. This is a failure to integrate an
existing capability, not evidence that the underlying defining sum or its
remainder bound is false.

\subsection{An existing public workaround}
The package's own arithmetic provides a useful control:
\begin{lstlisting}
s = AsymptoticAnalysis`AsymptoticExpansion[
  Zeta[x], {x, Infinity, Log[3]}, "Backend" -> "Package"];
AsymptoticAnalysis`SeriesAdd[s, -1]
(* GeneralizedSeries; Normal: (E^(-x))^Log[2] *)

s = AsymptoticAnalysis`AsymptoticExpansion[
  LerchPhi[1/2, 1, x], {x, Infinity, 3},
  "Backend" -> "Package"];
AsymptoticAnalysis`SeriesAdd[s, 1]
(* GeneralizedSeries; Normal: 1 - 2/x^2 + 2/x *)
\end{lstlisting}
The first normal expression equals $2^{-x}$ on the recorded real approach.
These controls demonstrate that the representation and arithmetic need not be
extended to handle the two translations. The missing step is reaching them
from the ordinary expression constructor.

The grouped arithmetic observation also emitted repeated native
\code{Power::infy} and \code{Greater::nord} messages. Their originating subcall
was not isolated. They are retained in the evidence record but are not elevated
to a separate defect or explained by a speculative mechanism.

\subsection{Why this is distinct from the earlier Lerch chart finding}
Report 24's N03 concerns an already constructed reciprocal-square Lerch chart:
recovering a signed coordinate such as $1/x^2$ can lose ordered arithmetic.
The current Lerch witness uses the direct argument $x$, and its existing
arithmetic control succeeds. The failure is earlier, at the construction of
\code{1 + LerchPhi[...]} as an ordinary expression. A signed quadratic chart
repair therefore does not address it.~\cite{report24}

Similarly, this is not another request for universal expression closure. A
rational affine wrapper is a narrow, explicitly checkable grammar whose
admitted examples require no new coefficient algebra. More general products,
multiple atoms, parameter-dependent factors, or nested functions should remain
separate admission decisions.

\subsection{A restricted lifting rule}
For a rational constant pair $b,c$, one varying admitted atom $A(x)$, and a
series result $S=e+\OO(R)$ for that atom, the desired transformation is simply
\begin{equation}
 c+bA(x)=c+be(x)+\OO(|b|R(x)).
\end{equation}
The package already implements the corresponding scalar multiplication and
addition. The prototype in \code{code/AffineDirichletExpansion.wl} performs
these steps, then applies the requested cutoff using the existing truncation
operation.

The prototype deliberately requires exactly one varying Zeta/Lerch atom,
degree-one dependence on that atom, and rational constant affine coefficients.
It declines quadratic dependence, inexact coefficients, variable-dependent
multipliers, and multiple distinct atoms. It adds one symbol under
\code{AsymptoticAudit`} and does not change any upstream definition or install
System upvalues.

\begin{lstlisting}
Get["AsymptoticAnalysis.wl"];
Get["AffineDirichletExpansion.wl"];

AsymptoticAudit`AffineDirichletExpansion[
  Zeta[x] - 1, {x, Infinity, Log[3]}]

AsymptoticAudit`AffineDirichletExpansion[
  1 + LerchPhi[1/2, 1, x], {x, Infinity, 3}]
\end{lstlisting}
The cutoff in this interface belongs to the special atom's coordinate. In
particular, a Zeta cutoff is not silently reinterpreted as a power of $1/x$.
The prototype promises an asymptotic remainder carried through the existing
arithmetic, not retention of every quantitative bound attached to the source
atom. It does not expose a term-goal option.

\subsection{The term-goal trap when integrating the rule upstream}
A direct implementation must count blocks of the final normalized expression,
not merely blocks requested from the atom. For example, if the user asks for
one term of $\zeta(x)-1$, asking the bare Zeta constructor for one term yields
only the constant one; subtracting one then erases the only retained block.
The desired leading nonzero block is $2^{-x}$.

For the rational affine grammar, the constant addition can cancel or introduce
at most one weight-zero block. A limited implementation can therefore request
one extra source block when a term goal is active, merge the affine constant,
and then apply the final goal. It must retain the primitive's honest remainder
when a cutoff prevents additional source information. This is a specific
one-block accounting rule for affine lifting, not a general backward-precision
planner.

The standalone prototype avoids this extra policy by accepting an explicit
cutoff only. The integrated constructor should eventually support both
constraints, but it should do so after N02's primitive request handling is fixed.

\section{N02: explicit cutoffs disable the term goal}
\begin{finding}
\textbf{New request-contract finding; successful baseline and candidate observations.}\par
\textbf{Locations:} the cutoff branch of \code{dirichletZetaForward} and the
stopping predicate in \code{dirichletLerchForward}.\par
\textbf{Effect:} extra returned blocks and, for Zeta, avoidable resource refusal.
The returned longer series is not mathematically false.
\end{finding}

\subsection{Contradictory request and result metadata}
The ordinary constructor already applies an integer term goal after the
exclusive cutoff: \code{makeForwardObject} takes at most the requested number
of retained blocks. The defining-sum constructors bypass that finishing path
and implement their own selection.~\cite{core,dirichlet}

The Lerch witness is:
\begin{lstlisting}
s = AsymptoticAnalysis`AsymptoticExpansion[
  LerchPhi[1/2, 1, x], {x, Infinity, 4},
  SeriesTermGoal -> 1, "Backend" -> "Package"];
{Normal[s], s["ReturnedTermCount"], s["RequestedTermGoal"]}
(* {6/x^3 - 2/x^2 + 2/x, 3, 1} *)
\end{lstlisting}
With the same cutoff and goal, the ordinary input
\code{1+1/x+1/x^2} returns only the constant one, records a returned count of
one, and places the first omitted power at $x^{-1}$. Thus the mismatch is
observable within the same package API, not a disputed inference from native
Wolfram option behavior.

For Zeta the corresponding observation is:
\begin{lstlisting}
s = AsymptoticAnalysis`AsymptoticExpansion[
  Zeta[x], {x, Infinity, Log[4]},
  SeriesTermGoal -> 1, "Backend" -> "Package"];
{Normal[s], s["ReturnedTermCount"], s["RequestedTermGoal"]}
(* {1 + 2^(-x) + 3^(-x), 3, 1} *)
\end{lstlisting}
The expectation is to retain one admitted nonzero block, here the constant one,
and to report a remainder beginning at $2^{-x}$. The public usage of
\code{SeriesTermGoal} as a term-count control is also consistent with the
native option's documented purpose, but the decisive evidence is the
repository's own ordinary constructor and recorded metadata.~\cite{wl-goal,core}

\subsection{The exact source mechanism}
Zeta distinguishes an automatic-cutoff request from an explicit-cutoff request.
In the automatic case it sets \code{count = goal}. In the explicit case it
binary-searches all integers whose logarithms lie below the cutoff and never
caps the result by the goal.

The Lerch path makes the exclusivity equally visible:
\begin{lstlisting}
If[If[cut === Automatic,
     Length[rows] >= goal,
     ! less[s + k, cut]],
   frontier = {s + k, coefficient, k}; Break[]];
\end{lstlisting}
The goal and cutoff are alternatives in that predicate. They should instead be
two independent reasons to stop retaining the next nonzero block. Validation
already accepts a positive integer goal together with an explicit exact cutoff,
so this is not an unsupported argument combination rejected by design.~\cite{dirichlet}

\subsection{A small term goal must constrain preflight work too}
A post hoc truncation alone is not enough for Zeta. The baseline also executes
\begin{lstlisting}
AsymptoticAnalysis`AsymptoticExpansion[
  Zeta[x], {x, Infinity, 1000000},
  SeriesTermGoal -> 1, "MaxTerms" -> 100,
  "Backend" -> "Package"]
(* Failure["ResourceLimit", ...] *)
\end{lstlisting}
Its resource message says that the exponential-coordinate cutoff exceeds
\code{MaxTerms}. Under the combined request semantics, however, the only
retained block is one, and the first omitted integer is two. No enumeration
near $e^{1000000}$, nor even a 100-term expansion, is needed.

The candidate therefore restricts the binary-search interval itself. When the
goal is a positive integer $g$, the upper search sentinel becomes
$\min(g,L)+1$, where $L$ is the existing term budget. A cutoff-based resource
refusal is necessary only when the cutoff would require more than $L$ terms
\emph{and} no active goal reduces the retained count to at most $L$.

\subsection{The three-edit candidate}
The staged patch changes exactly these expressions:
\begin{lstlisting}
(* Zeta preflight: old condition *)
less[Log[limit + 1], cut]

(* New condition *)
less[Log[limit + 1], cut] &&
  (goal === Automatic || goal > limit)

(* Zeta search interval *)
low = 0;
high = If[goal === Automatic,
  limit + 1, Min[limit, goal] + 1];

(* Lerch stopping rule, after a nonzero coefficient is found *)
(cut =!= Automatic && ! less[s + k, cut]) ||
  (goal =!= Automatic && Length[rows] >= goal)
\end{lstlisting}
All coefficient formulas, existing domain tests, and bound constructors remain
unchanged. In particular, the Lerch loop still finds a \emph{nonzero} frontier
coefficient before stopping; it does not count vanished moment coefficients
as returned terms.

\subsection{Why the selected frontier remains correct}
For Zeta, let $m(h)$ be the number of positive integers with $\log n<h$.
With both constraints active, the retained count is
\begin{equation}
 N=\min\{m(h),g\}.
\end{equation}
The new binary search computes this minimum without needing $m(h)$ when it is
larger than the goal. The first omitted integer is still $N+1$, so the existing
remainder constructor can use $\rho=\log(N+1)$ and its existing tail metadata.
The proof does not depend on approximating $e^h$ numerically. Exact comparisons
continue to decide logarithmic boundary equality.

For Lerch, suppose the nonzero coefficient indices are
$k_0<k_1<\cdots$. The loop retains the first indices whose weights $s+k_j$ lie
below the cutoff, stopping also when $g$ such blocks have been retained. The
next nonzero coefficient supplies the same kind of frontier as before. Only
the selected frontier changes; the existing Taylor-moment bound is evaluated
at that frontier's index.

\subsection{Observed candidate results}
The three anchor strings were each verified to occur exactly once in the
imported pinned standalone. After replacement, that temporary file was loaded
in a fresh Wolfram evaluator. The following outputs were observed, with
\code{MaxTerms -> 100}:
\begin{longtable}{@{}P{0.24\textwidth}P{0.16\textwidth}P{0.12\textwidth}P{0.39\textwidth}@{}}
\toprule
Input & Cutoff & Goal & Candidate normal expression and remainder\\\midrule\endhead
$\zeta(x)$ & $\log4$ & 1 & $1+\OO(2^{-x})$; returned count 1.\\
$\Phi(1/2,1,x)$ & 4 & 1 & $2/x+\OO(x^{-2})$; returned count 1.\\
$\zeta(x)$ & $10^6$ & 1 & $1+\OO(2^{-x})$; no resource refusal.\\
$\zeta(x)$ & $\log4$ & Automatic & $1+2^{-x}+3^{-x}+\OO(4^{-x})$; unchanged.\\
$\Phi(1/2,1,x)$ & 4 & Automatic & $2/x-2/x^2+6/x^3+\OO(x^{-4})$; unchanged.\\
\bottomrule
\end{longtable}

The candidate changes achieved precision when it honors the smaller term goal;
it does not claim that a one-term result has the longer expansion's remainder.
That distinction is essential. Merely deleting two terms while retaining the
old fourth-order error would introduce a mathematical error that the baseline
does not have.

\subsection{Novelty and integration boundary}
Earlier reports discuss configured backend defaults and special inverse power
cutoffs. N02 is not a repeat of those dispatch/default bugs: the options in
these witnesses are explicit, already parsed, and accurately stored in the
result. The two defining-sum loops then fail to apply them together. The patch
is local to those loops rather than another general request-normalization
proposal.~\cite{wave3,wave5}

The Python staging utility checks all three anchors before writing anything,
refuses to overwrite existing output, and writes a canonical module plus a
unified diff into a separate directory. It does not modify the checkout in
place. The successful native experiment edited the generated standalone's
corresponding source text; a canonical-module edit followed by the upstream
builder remains a separate integration step to execute before adoption.

\section{A narrow development option: integer-indexed Dirichlet jets}
\subsection{The representation opportunity}
The Zeta adapter currently exposes the sequence $\log1,\log2,\log3,\ldots$ as
ordinary exact real weights. That is mathematically appropriate for its
positive coordinate. But an algebra dedicated to this defining-sum family can
retain the \emph{integer index} before converting to logarithmic weights:
\begin{equation}\label{eq:dirichlet}
 A(S)=\sum_{n\geq1}a_n n^{-S}.
\end{equation}
For positive integer indices, order and multiplication reduce to
\begin{equation}
 \log n<\log m\iff n<m,
 \qquad \log n+\log m=\log(nm).
\end{equation}
Thus this family has a more specific algebra than arbitrary real exponent
weights. Its coefficients can be combined using integer multiplication and
divisor relations. This is not a claim that the current generic jet algebra
cannot express the same functions, nor a measured claim that it is slow.
It is a concrete optional specialization.

The supplied independent prototype restricts coefficients to rational numbers
and the growth exponent to a nonnegative integer. Those restrictions make
coefficient operations, majorant selection, and tail evaluation at integer $S$
exact using only the Python standard library. A future Wolfram implementation
could use a wider exact coefficient ring without changing the indexing idea.

\subsection{The prefix invariant}
A jet records a positive integer $N$, all coefficients for $1\leq n\leq N$,
and a global coefficient bound
\begin{equation}\label{eq:growth}
 |a_n|\leq Cn^p\quad(n\geq1),\qquad C\geq0,
 \quad p\in\mathbb Z_{\geq0}.
\end{equation}
The coefficient dictionary is sparse, but its \emph{known prefix is dense
information}: an absent dictionary entry inside $1,\ldots,N$ means a known zero.
An absent entry beyond $N$ is unknown, not zero. The accessor refuses such an
out-of-prefix query.

The public constructor declares the infinite bound and checks it on every
supplied prefix coefficient. That finite check cannot establish an arbitrary
infinite assertion. The Zeta factory establishes its bound by construction
because every coefficient is one. The finite-polynomial factory likewise
knows its entire sequence. Arithmetic methods then transport a supplied bound
by the propositions below. The distinction is part of the class contract, not
hidden behind the word ``certified.''

The prototype's retained expression is
\begin{equation}
 A_N(S)=\sum_{n=1}^{N}a_n n^{-S}.
\end{equation}
Using a prefix rather than a count of nonzero coefficients makes cancellation
and sparse storage unambiguous. A user-facing term goal can be implemented
above this layer, but it must not redefine unknown prefix entries as zeros.

\subsection{An exact elementary tail bound}
\begin{proposition}[Tail from a coefficient majorant]\label{prop:tail}
Assume \eqref{eq:growth}, let $M=N+1$, and let $S>p+1$. Then the Dirichlet series
converges absolutely and
\begin{equation}\label{eq:tail}
 |A(S)-A_N(S)|
 \leq C M^{p-S}\left(1+\frac{M}{S-p-1}\right).
\end{equation}
\end{proposition}
\begin{proof}
Put $t=S-p>1$. Absolute values reduce the omitted tail to
$C\sum_{n=M}^{\infty}n^{-t}$. The decreasing summand is bounded by its first
term plus the integral from $M$ to infinity. Evaluating that integral gives
$CM^{-t}(1+M/(t-1))$, which is \eqref{eq:tail}.
\end{proof}

This is the same elementary comparison principle used by the existing Zeta
adapter, now applied to a transported coefficient-growth bound. The incremental
point is the closed coefficient-bound calculus below, not the integral test
itself. At integer $S>p+1$, every quantity returned by the prototype's
\code{tail_bound} method is rational and computed exactly.

\subsection{Addition and scalar multiplication}
For two coefficient sequences with bounds $(C_a,p_a)$ and $(C_b,p_b)$,
addition admits
\begin{equation}
 C_{a+b}=C_a+C_b,\qquad p_{a+b}=\max(p_a,p_b).
\end{equation}
Scalar multiplication by a rational $r$ admits $(|r|C_a,p_a)$. The known prefix
of a binary operation is the intersection of the input prefixes. These rules
are conservative when operands are correlated: canceling retained coefficients
does not, by itself, cancel independently declared unknown tails.

\subsection{Multiplication is Dirichlet convolution}
\begin{proposition}[Coefficient and majorant transport]\label{prop:product}
In a common half-line of absolute convergence, the product of two series
\eqref{eq:dirichlet} has coefficients
\begin{equation}\label{eq:convolution}
 c_n=\sum_{d\mid n}a_d b_{n/d}.
\end{equation}
If the input sequences obey their declared bounds, a valid elementary output
bound is
\begin{equation}
 |c_n|\leq C_aC_b n^{\max(p_a,p_b)+1}.
\end{equation}
All output coefficients through $N=\min(N_a,N_b)$ are determined by the two
input prefixes.
\end{proposition}
\begin{proof}
Absolute convergence permits regrouping the product by $n=dm$, which gives
\eqref{eq:convolution}. Every divisor factor satisfies
$d^{p_a}(n/d)^{p_b}\leq n^{\max(p_a,p_b)}$. There are at most $n$ positive
divisors of $n$, giving the stated bound. For $n\leq N$, every divisor and
complementary divisor is at most $N$, so no unknown coefficient is used.
\end{proof}

The divisor-count estimate is intentionally loose. It keeps $p$ integral and
all subsequent majorants rational at integer $S$. A sharper estimate such as
$\tau(n)\leq2\sqrt n$ would reduce the growth exponent at the price of a
larger exponent domain in the implementation. Nothing in the prototype claims
an optimal bound or an optimal abscissa of convergence.

For dense known prefixes, the retained pair count is exactly
\begin{equation}\label{eq:pairs}
 H_N=\sum_{d=1}^{N}\left\lfloor\frac Nd\right\rfloor.
\end{equation}
The elementary harmonic-sum comparison gives $H_N\leq N(1+\log N)$.
A row-major loop therefore stops its inner traversal at $m\leq N/d$ instead of
forming an $N$ by $N$ product and discarding most pairs afterwards. Sparse input
support reduces the actual work further.

\subsection{Reciprocals with a coefficient bound, not only a formal recurrence}
A reciprocal prefix can be computed whenever $a_1\ne0$:
\begin{equation}\label{eq:recurrence}
 b_1=\frac1{a_1},\qquad
 b_n=-\frac1{a_1}\sum_{\substack{d\mid n\\d\geq2}}a_d b_{n/d}
 \quad(n>1).
\end{equation}
The missing ingredient in a purely formal implementation would be a useful
global bound for the new sequence. The following elementary construction
supplies one.

\begin{proposition}[A computable reciprocal majorant]\label{prop:reciprocal}
Suppose $a_1\ne0$ and \eqref{eq:growth} holds. Choose an integer $q>p+1$ for which
\begin{equation}\label{eq:kappa}
 \kappa=\frac C{|a_1|}\,2^{p-q}
 \left(1+\frac2{q-p-1}\right)<1.
\end{equation}
The sequence defined by \eqref{eq:recurrence} satisfies
\begin{equation}
 |b_n|\leq\frac{n^q}{|a_1|}\quad(n\geq1).
\end{equation}
Its absolutely convergent Dirichlet series represents $1/A(S)$ for $S>q+1$.
\end{proposition}
\begin{proof}
Set $D=1/|a_1|$. The base coefficient satisfies $|b_1|=D$. If the bound holds
for all indices below $n$, then each $n/d$ in \eqref{eq:recurrence} is smaller
than $n$, and hence
\begin{align*}
 |b_n|&\leq Dn^q\frac1{|a_1|}
       \sum_{\substack{d\mid n\\d\geq2}}|a_d|d^{-q}\\
 &\leq Dn^q\frac C{|a_1|}\sum_{d=2}^{\infty}d^{p-q}
 \leq Dn^q\kappa < Dn^q.
\end{align*}
The final bound on the sum is the first-term-plus-integral estimate beginning
at two. This proves the coefficient inequality by strong induction. For
$S>q+1$, both the original and reciprocal series converge absolutely. Their
product has coefficient one at index one and zero at every larger index by
\eqref{eq:recurrence}. Thus their product is one, proving the reciprocal claim.
\end{proof}

The expression in \eqref{eq:kappa} tends to zero as $q$ increases, so a suitable
integer exists. With rational $C$ and $a_1$, the search can compare $\kappa$ to
one using exact fractions. The implementation bounds both this search and the
coefficient work. Budget exhaustion raises an explicit exception instead of
returning an unproved reciprocal bound.

A sieve-style evaluation avoids asking for divisors separately at every index.
After coefficient $b_m$ is known, its contributions $a_d b_m$ are accumulated
at indices $dm\leq N$, $d\geq2$. All contributions required for the next
coefficient have arrived from smaller indices. This is the same retained-pair
geometry as \eqref{eq:pairs}.

\subsection{Executable examples}
\begin{lstlisting}[language=Python]
from dirichlet_jet import DirichletJet

z = DirichletJet.zeta(128)
z2 = z.multiply(z)
reciprocal = z.reciprocal()
identity = z.multiply(reciprocal)

assert dict(identity.coefficients) == {1: 1}
assert z2.coefficient(12) == 6
assert reciprocal.coefficient(6) == 1
assert reciprocal.coefficient(12) == 0

finite_value = reciprocal.evaluate_prefix(10)
absolute_tail = reciprocal.tail_bound(10)  # exact Fraction
\end{lstlisting}
For $\zeta^2$, the coefficients count the divisors of $n$. For the reciprocal
of Zeta, the computed prefix agrees with an independent elementary
prime-factorization implementation of the M\"obius function. These are useful
structural oracles, not comparisons with a second transcription of the same
convolution loop.

The example does not assert that the conservative reciprocal majorant is the
sharp bound for M\"obius coefficients. A source-aware specialized factory could
provide a stronger bound, but the implemented reciprocal operation intentionally
uses only its input's declared coefficient majorant.

\section{Performance evidence and implementation tradeoffs}
\subsection{What was actually measured}
The independent Python convolution prototype was measured on dense Zeta
prefixes. Input construction was outside each timed interval. Each size was
repeated five times; the table reports median elapsed seconds in this session.
These are not Wolfram or Mathics timings, and no speedup over the existing
package has been measured.

\begin{center}
\begin{tabular}{rrrr}
\toprule
Known prefix $N$ & Retained pairs $H_N$ & Unfiltered $N^2$ pairs & Median seconds\\\midrule
64 & 280 & 4,096 & 0.00047\\
256 & 1,466 & 65,536 & 0.00216\\
1,024 & 7,262 & 1,048,576 & 0.01086\\
4,096 & 34,720 & 16,777,216 & 0.04985\\\bottomrule
\end{tabular}
\end{center}

The exact pair counts explain the loop's search geometry independently of the
timer. The implementation still pays for fraction arithmetic, dictionary
updates, output normalization, and prefix-bound checks. The measurement does
not estimate peak memory, simplifier costs in Wolfram, or the behavior of a
symbolic coefficient ring. The complete platform string and individual
summary values are in \code{evidence/python-benchmark.json}; the benchmark
script is included.

\subsection{Why integer indexing is useful even without a speed claim}
Integer indexing removes a particular proof obligation from this family:
recognizing equality or order between sums of weights that are known to be
logarithms of integer products. A coefficient of $6^{-S}$ produced from
$2^{-S}3^{-S}$ has index six immediately. There is no need to ask a generic
symbolic simplifier to recognize $\log2+\log3=\log6$.

This observation justifies a data representation, not a blanket recommendation
to replace the existing real-weight implementation. Irrational-power germs,
Fourier-polynomial coefficients, and the package's other scales require their
own contracts. The integer representation should expose a faithful conversion
back to the existing $w^{\log n}$ rows rather than silently assuming all real
weights have integer logarithmic indices.

\subsection{What the prototype intentionally does not implement}
The Python class does not integrate with \code{GeneralizedSeries}; it is an
independent reference model. It does not implement arbitrary real or complex
coefficients, noninteger phase evaluation, general phase composition,
parameter-uniform estimates, automatic source recognition, or exact
cancellation of correlated unknown tails. Its finite-polynomial factory also
uses the generic conservative tail interface, rather than exploiting finite
support to report a zero tail through a second exactness channel.

Those are explicit prototype boundaries. They are not additional defects
attributed to the upstream repository. The implementation is small enough to
serve as an oracle for a future family-specific module and to test whether
integer indexing is worth integrating before complicating the public result
schema.

\section{Wolfram 15 and Mathics: a targeted compatibility assessment}
\subsection{What the successful Wolfram run establishes}
The actual native kernel identified itself as
\begin{lstlisting}
15.0.0 for Linux x86 (64-bit) (May 6, 2026)
\end{lstlisting}
It loaded the complete imported standalone and executed the public witness
calls. The term-goal candidate also executed there. Therefore N01 and N02 are
not merely predictions based on a different Wolfram version, a Mathics
emulation, or an independent Python reimplementation.

The no-goal controls are especially important for N02: the candidate must not
change the series selected by the cutoff alone. The observed unchanged
Zeta and Lerch expressions demonstrate that limited property for the two
chosen inputs. They do not establish every parameter case or an end-to-end
release gate.

\subsection{What source inspection says about Mathics}
The repository isolates several host differences behind Mathics-specific
helpers: assumption/sign reasoning, simplification, Taylor expansion,
numerical evaluation, inverse-branch reasoning, and compatibility shims.
The source contains explicit guards around native behaviors that the package
does not treat as interchangeable with Wolfram. This architecture matters when
assessing a proposed edit: a fix should use the package's existing exact-order
and validation helpers rather than add an unreviewed dependency on a native
symbolic oracle.~\cite{mathics-assumptions,mathics-simplification,mathics-compat}

The term-goal patch changes only Boolean conditions, integer bounds, and the
existing \code{less} comparison calls. It does not introduce a new special
function, a native-only calculus operation, or a new root finder. Both affected
constructors are shared canonical source, so the control-flow issue is not
inherently Wolfram-specific. Whether every corresponding public witness reaches
that path and produces the same output under a particular Mathics version
still requires an actual run.

The affine prototype similarly reuses the package's public arithmetic, but its
complete loader/option/pattern behavior has not been validated under Mathics.
Its code is intended to be straightforward to test, not presented as a
compatibility-certified drop-in. Mathics is an independent implementation of
the Wolfram Language, so successful Wolfram execution cannot substitute for
that second run.~\cite{mathics-site}

\subsection{A focused cross-host acceptance set}
For the proposed changes, the useful acceptance cases are narrowly defined.
The bare-atom controls must still succeed; the combined goal/cutoff cases must
retain the same selected blocks in both hosts; a no-goal request must remain
unchanged; and the large-cutoff, one-term Zeta request must no longer fail solely
because of the inactive work implied by the cutoff.

Boundary cases should additionally cover an exact Zeta cutoff at $\log m$,
a goal larger than the number of available finite Lerch blocks, an invalid
goal, a goal larger than the budget but an earlier cutoff, and an affine
translation whose weight-zero block cancels. These cases follow from the
specific changes in this article. They are not a restatement of the existing
general cross-kernel testing program.

The packaged \code{RunProbes.wl} uses plain records rather than relying on a
shared implementation of \code{TestReport}. It performs no download or
repository mutation. It is a proposed reproduction driver; the successful
observations in this review came from separately issued focused calls, not
from execution of the complete packaged file.

\section{Concrete integration and development sequence}
\subsection{First: land the bounded request fix independently}
Stage the three-edit canonical change, rebuild the standalone using the
repository's own generation path, and run the focused N02 comparisons on both
hosts. Keep this change independent of affine lifting and the integer-indexed
prototype. Its coefficient formulas and result schema need not change.

The acceptance condition is not merely that the normal expressions look
shorter. Returned counts, first omitted terms, remainder powers, and explicit
bounds must all correspond to the selected shorter prefix. The native candidate
already demonstrates the expected frontier change for the selected examples.

\subsection{Second: admit a small affine grammar at the constructor boundary}
Use the rational-affine prototype as a behavioral sketch, not as an instruction
to call a public constructor recursively on every subtree. The integrated path
should identify one admitted atom, prove the permitted coefficient conditions,
construct that atom in its own coordinate, apply the exact affine operations,
and finish the request once on the normalized result.

Its diagnostic should distinguish an unsupported \emph{wrapper grammar} from
an unsupported special function. For example, a rejected variable-dependent
multiplier should say that the affine defining-sum route requires constant
coefficients. A later native fallback may still be attempted under the user's
chosen backend policy, but its incidental coefficient error should not be the
only explanation available when the special leaf itself is recognized.

\subsection{Third: evaluate integer-indexed algebra as a separate feature}
A small initial Wolfram-facing interface could construct an integer-indexed
Dirichlet prefix, combine prefixes by addition and convolution, compute an
admitted reciprocal, and export $\{\log n,a_n\}$ rows plus a transported
majorant. It should first support a single common phase $S$ and rational
coefficients, matching the independent oracle.

Only after that core agrees with the independent tests should it widen the
coefficient ring. Differentiation in $S$ is a particularly natural next
extension because
\begin{equation}
 \frac{d^k}{dS^k}n^{-S}=(-\log n)^k n^{-S}.
\end{equation}
A simple bound $(\log n)^k\leq n^k$ for $n\geq1$ would preserve an elementary
integer growth exponent, though it is conservative. For Zeta, the corresponding
termwise derivative formulas in the absolute-convergence half-plane are also
recorded in DLMF.~\cite{dlmf-zeta}

This proposed extension is tied to the integer-indexed family and its
coefficient-majorant proof. It is not a proposal to claim differentiability
from an arbitrary asymptotic remainder. The prototype currently implements
neither derivative coefficients nor that extended coefficient ring.

A second controlled extension is an integer affine phase shift:
$A(S+b)=\sum a_n n^{-b}n^{-S}$. For integer $b$, rational coefficients remain
rational; for negative $b$, the coefficient growth exponent increases by $-b$.
This can be added without admitting unrelated exponential rates. Arbitrary
phase compositions should remain outside this initial feature.

\subsection{What should not be coupled to these changes}
There is no need to change the entire result schema, replace the general
precision calculus, redesign all native dispatch, or promise arbitrary
transseries closure to fix N01 and N02. Those wider topics already have their
own record. A family-specific implementation should earn its place by reducing
an observed failure and agreeing with a small independent algebra, rather
than becoming a prerequisite for unrelated repairs.

\section{Artifacts, reproduction, and final assessment}
\subsection{Archive contents}
\begin{longtable}{@{}P{0.46\textwidth}P{0.47\textwidth}@{}}
\toprule
Artifact & Purpose and evidence status\\\midrule\endhead
\code{article/review.tex}, \code{article/review.pdf} & Source and rendered article.\\
\code{code/patch_dirichlet_goals.py} & Guarded staging utility for the three-edit candidate; synthetic anchor tests passed, and the corresponding edits passed five native controls.\\
\code{code/AffineDirichletExpansion.wl} & Restricted additive/scalar lifting prototype; complete helper execution remains unverified.\\
\code{code/RunProbes.wl} & Local-file-only Wolfram/Mathics reproduction driver; full packaged driver not executed here.\\
\code{code/dirichlet_jet.py} & Exact-rational integer-indexed prefix, convolution, reciprocal, and coefficient-bound implementation.\\
\code{code/benchmark_dirichlet.py} & Reproducible benchmark of that independent Python prototype only.\\
\code{tests/test_dirichlet_jet.py} & 32 passing independent test methods, including four synthetic patch-fixture methods.\\
\code{evidence/} & Native transcriptions, source coverage, novelty ledger, independent status, complete Python test output, and prototype benchmark results.\\
\bottomrule
\end{longtable}

\subsection{Running the independent checks}
From the archive root:
\begin{lstlisting}
python -m unittest discover -s tests -v
python code/benchmark_dirichlet.py
\end{lstlisting}
The tests use only the standard library. They check divisor-count coefficients,
a prime-factorization M\"obius oracle, reciprocal identities, random exact
convolutions, prefix boundaries, rational majorants, input validation,
immutability, and explicit work budgets. Finite partial-tail comparisons are
sanity checks accompanying the displayed infinite-tail proofs, not substitutes
for those proofs.

To stage the candidate without altering the source checkout:
\begin{lstlisting}
python code/patch_dirichlet_goals.py /path/to/Asymptotic /path/to/stage
\end{lstlisting}
The output directory must not contain the staged artifacts already. A changed
or already-patched source is rejected by the anchor checks. Apply the resulting
canonical change in a disposable checkout and use the upstream builder before
comparing the generated standalone. This review does not supply a replacement
copy of the entire upstream repository.

For a fresh Wolfram or Mathics session, set a local path before loading the
probe driver:
\begin{lstlisting}
AuditPackagePath = "/path/to/Asymptotic/AsymptoticAnalysis.wl";
AuditOutputPath = "/path/to/observations.wl";
Get["/path/to/audit/code/RunProbes.wl"];
\end{lstlisting}
Use a separate fresh process for the baseline and candidate. The optional
output path is a plain-text observation record, not an acceptance receipt.

\subsection{Assessment}
The new actionable defects are small and localized: a missing affine admission
route, and a conditional that treats two simultaneous constraints as alternatives.
Their importance is not that they undermine the package's entire mathematics.
They make capabilities already present in the repository unexpectedly hard to
use and make request behavior depend on the selected special-function path.

The term-goal candidate is the most immediately supported change: it was
executed against the pinned standalone and preserves the selected no-goal
controls. The affine helper is a restricted implementation sketch with
successful underlying arithmetic controls, not a fully validated integration.
The integer-indexed Dirichlet prototype is a separate, implemented development
option with a precise algebra and explicit coefficient-bound assumptions.

This is the full incremental conclusion supported by the source inspection
and executions obtained here. It does not manufacture additional findings from
already catalogued issues, transport failures, or features outside an explicitly
admitted mathematical contract.

\begin{thebibliography}{99}
\small\setlength{\itemsep}{2pt}\setlength{\parskip}{0pt}
\bibitem{reviewindex} Asymptotic Contributors. \emph{External code-review index and retirement records}, pinned revision. \src{external-reports/code-review/README.md}{Repository review index}.
\bibitem{register} Asymptotic Contributors. \emph{Code review status and implementation register}, pinned revision. \src{docs/development/CODE_REVIEW_STATUS.md}{Maintained register}.
\bibitem{wave3} Asymptotic Contributors. \emph{Wave-3 review index}. \src{external-reports/code-review/wave-3/README.md}{Wave-3 packages and finding summaries}.
\bibitem{wave5} Asymptotic Contributors. \emph{Wave-5 review index}. \src{external-reports/code-review/wave-5/README.md}{Wave-5 packages and finding summaries}.
\bibitem{report24} \emph{Incremental technical audit}, retained report 24, especially the reciprocal-square Lerch discussion. \src{external-reports/code-review/wave-3/code-review-24/article/audit.tex}{Original article source}.
\bibitem{core} Asymptotic Contributors. \emph{AsymptoticAnalysis.wl}, canonical kernel, especially \code{localCoordinate}, \code{forwardCore}, and \code{makeForwardObject}. \src{src/Kernel/AsymptoticAnalysis.wl}{Pinned canonical source}; relevant request/object assembly begins near lines 610--740.
\bibitem{dirichlet} Asymptotic Contributors. \emph{DirichletSpecialFunctions.wl}, especially \code{dirichletSpecialMake}, \code{dirichletZetaForward}, \code{dirichletLerchForward}, and the final whole-expression dispatcher. \src{src/Kernel/DirichletSpecialFunctions.wl}{Pinned defining-sum module}.
\bibitem{arithmetic} Asymptotic Contributors. \emph{SeriesArithmetic.wl}, regular operands and held arithmetic dispatch. \src{src/Kernel/SeriesArithmetic.wl}{Pinned source}, especially the first 180 lines.
\bibitem{operations} Asymptotic Contributors. \emph{SeriesOperations.wl}, ordinary series data, operations, and result construction. \src{src/Kernel/SeriesOperations.wl}{Pinned source}.
\bibitem{envelope} Asymptotic Contributors. \emph{SeriesEnvelopeArithmetic.wl}, compatible approaches and conservative composite arithmetic. \src{src/Kernel/SeriesEnvelopeArithmetic.wl}{Pinned source}.
\bibitem{mathics-assumptions} Asymptotic Contributors. \emph{MathicsAssumptions.wl}, finite-real and sign consequence helpers. \src{src/Kernel/MathicsAssumptions.wl}{Pinned source}.
\bibitem{mathics-simplification} Asymptotic Contributors. \emph{MathicsSimplification.wl}, host simplification and limit adapters. \src{src/Kernel/MathicsSimplification.wl}{Pinned source}.
\bibitem{mathics-compat} Asymptotic Contributors. \emph{MathicsCompatibility.wl}, bounded compatibility helpers. \src{src/Kernel/MathicsCompatibility.wl}{Pinned source}.
\bibitem{mathics-site} Mathics3 Project. \emph{Mathics3}. Official project site, consulted 10 September 2026. \url{https://mathics.org/}.
\bibitem{wl-goal} Wolfram Research. \emph{SeriesTermGoal}. Wolfram Language documentation, consulted 10 September 2026. \url{https://reference.wolfram.com/language/ref/SeriesTermGoal.html}.
\bibitem{dlmf-zeta} NIST Digital Library of Mathematical Functions. \emph{Section 25.2: Definition and expansions}, especially equations 25.2.1 and 25.2.7. \url{https://dlmf.nist.gov/25.2}.
\bibitem{dlmf-lerch} NIST Digital Library of Mathematical Functions. \emph{Section 25.14: Lerch's transcendent}, especially its defining series. \url{https://dlmf.nist.gov/25.14}.
\end{thebibliography}
\end{document}
